\documentclass[11pt]{article}
\usepackage[a4paper,margin=2.1cm]{geometry}
\usepackage{ctex}
\usepackage{amsmath,amssymb}
\usepackage{xcolor}
\usepackage{tikz}
\usepackage{tcolorbox}
\usepackage{array}
\usepackage{booktabs}
\usetikzlibrary{arrows.meta,calc,positioning}
\tcbuselibrary{skins,breakable}

\definecolor{curve}{RGB}{30,90,180}
\definecolor{secline}{RGB}{220,60,60}
\definecolor{tang}{RGB}{230,140,20}
\definecolor{auxc}{RGB}{20,150,90}
\definecolor{grayt}{RGB}{110,110,110}
\definecolor{boxbg}{RGB}{245,249,255}

\title{\textbf{二阶导符号与凸性}\\[2pt]
\large 由 $f''(x)\ge 0$ 证明 $f$ 是下凸函数（充分性），及图像解释}
\author{}
\date{}

\begin{document}
\maketitle
\vspace{-6mm}

\begin{tcolorbox}[colback=boxbg,colframe=curve!70,title={\bfseries 核心命题}]
设 $f$ 在区间 $I$ 上二阶可导。若对一切 $x\in I$ 有 $\;f''(x)\ge 0$，则 $f$ 在 $I$ 上是\emph{下凸函数}。\\
等价结论：$f'$ 在 $I$ 上单调增加，且曲线恒在其任意点切线之上。
\end{tcolorbox}

\textbf{要证．}对任意 $x_1,x_2\in I$ 及 $\lambda\in(0,1)$，
\[
f\big(\lambda x_1+(1-\lambda)x_2\big)\ \le\ \lambda f(x_1)+(1-\lambda)f(x_2). \tag{$\ast$}
\]

% =====================================================================
\section{证明（逐行解读）}

\textbf{第 1 步：$f''\ge 0\Rightarrow f'$ 单调增加.}
因为 $f''=(f')'\ge 0$，由「导数非负 $\Rightarrow$ 函数单调不减」得 $f'$ 在 $I$ 上单调增加。
（这是全证明的核心引擎。）

\vspace{2mm}
\textbf{第 2 步：取内点 $x_0=\lambda x_1+(1-\lambda)x_2$，不妨设 $x_1<x_2$.}
则 $x_1<x_0<x_2$，且把「$x_0$ 的位置」换算成水平增量：
\[
x_1-x_0=(1-\lambda)(x_1-x_2),\qquad x_2-x_0=\lambda(x_2-x_1). \tag{1}
\]

\vspace{2mm}
\textbf{第 3 步：在 $[x_1,x_0]$ 与 $[x_0,x_2]$ 上分别用拉格朗日中值定理.}
存在 $\eta_1\in(x_1,x_0)$、$\eta_2\in(x_0,x_2)$，使
\[
f(x_1)=f(x_0)+f'(\eta_1)(x_1-x_0),\qquad
f(x_2)=f(x_0)+f'(\eta_2)(x_2-x_0). \tag{2}
\]

\vspace{2mm}
\textbf{第 4 步：关键放缩——用 $f'$ 单调增加.}
由 $x_1<\eta_1<x_0<x_2$ 且 $x_0<\eta_2<x_2$，有 $\eta_1<x_0<\eta_2$，故
\[
f'(\eta_1)\ \le\ f'(x_0)\ \le\ f'(\eta_2).
\]
再注意 $x_1-x_0<0$（$x_1<x_0$）、$x_2-x_0>0$，代入 (2) 并乘上对应符号，得
\[
f(x_1)\ \ge\ f(x_0)+f'(x_0)(x_1-x_0),\qquad
f(x_2)\ \ge\ f(x_0)+f'(x_0)(x_2-x_0). \tag{3}
\]
（(3) 的几何意义正是：\emph{曲线在过 $(x_0,f(x_0))$ 的切线上方}。）

\vspace{2mm}
\textbf{第 5 步：分别乘 $\lambda$ 与 $1-\lambda$ 再相加.}
将 (3) 两式分别乘 $\lambda$ 与 $1-\lambda$ 后相加：
\[
\lambda f(x_1)+(1-\lambda)f(x_2)
\ \ge\ f(x_0)+\Big[\lambda(x_1-x_0)+(1-\lambda)(x_2-x_0)\Big]f'(x_0).
\]
方括号内恰为 $0$——因为 $x_0=\lambda x_1+(1-\lambda)x_2$ 正是加权平均，故
$\lambda(x_1-x_0)+(1-\lambda)(x_2-x_0)=\lambda x_1+(1-\lambda)x_2-x_0=0$。
于是
\[
\lambda f(x_1)+(1-\lambda)f(x_2)\ \ge\ f(x_0)=f\big(\lambda x_1+(1-\lambda)x_2\big).
\]
这正是 $(\ast)$，即 $f$ 为下凸函数。$\qquad\blacksquare$

\begin{tcolorbox}[colback=boxbg,colframe=auxc!70,title={\bfseries 逻辑链小结}]
$f''\ge0 \;\Rightarrow\; f'$ 单调增 $\;\Rightarrow\;$ 两段平均斜率 $f'(\eta_1)\le f'(x_0)\le f'(\eta_2)$
$\;\Rightarrow\;$ 曲线在切线上方 (3) $\;\Rightarrow\;$ 加权相加 $\Rightarrow$ 凸性定义 $(\ast)$。
\end{tcolorbox}

% =====================================================================
\section{图像解释}

\subsection*{图 1：曲线在切线上方（第 4 步的几何本质）}
$f''\ge0$ 表示曲线「向上弯」。于是过任一点作的切线，整条曲线都在它\emph{上方}。
第 4 步的 (3) 式就是这张图：
\[
f(x_1)\ge f(x_0)+f'(x_0)(x_1-x_0),\quad f(x_2)\ge f(x_0)+f'(x_0)(x_2-x_0).
\]

\subsection*{图 1：切线支撑 + 弦在曲线上方}

\begin{center}
\begin{tikzpicture}[scale=1.08,>=Stealth]
  \draw[->,gray] (0,0)--(6.8,0) node[right,black]{$x$};
  \draw[->,gray] (0,0)--(0,5.4) node[above,black]{$f$};
  % f(x)=0.7+0.20*(x-1.2)^2
  % x1=1.4 -> 0.7+0.20*0.04=0.708 ; x0=3.4 -> 0.7+0.20*4.84=1.668 ; x2=5.6 ->0.7+0.20*19.36=4.572
  % f'(x)=0.40*(x-1.2) ; x0=3.4 -> 0.40*2.2=0.88
  \draw[curve,very thick,domain=0.7:6.0,samples=100,smooth]
       plot (\x,{0.7+0.20*(\x-1.2)*(\x-1.2)});
  % 切线 at x0=3.4 slope0.88
  \draw[tang,very thick] (0.9,{1.668+0.88*(0.9-3.4)})--(6.2,{1.668+0.88*(6.2-3.4)});
  \node[tang,font=\scriptsize,below right] at (5.6,{1.668+0.88*(5.6-3.4)}){切线（斜率 $f'(x_0)$）};
  % 弦
  \draw[secline,very thick] (1.4,0.708)--(5.6,4.572);
  \node[secline,font=\scriptsize,below right] at (4.9,3.4){弦};
  \fill (1.4,0.708) circle (2pt) node[below left,black]{$x_1$};
  \fill (3.4,1.668) circle (2pt) node[below,black]{$x_0$};
  \fill (5.6,4.572) circle (2pt) node[above left,black]{$x_2$};
  \draw[dashed,grayt] (3.4,1.668)--(3.4,0);
  % 切线在 x1 处的外推点
  \fill[tang] (1.4,{1.668+0.88*(1.4-3.4)}) circle (1.6pt);
  \node[tang,font=\scriptsize,below left] at (1.4,{1.668+0.88*(1.4-3.4)}){切线外推值 $\le f(x_1)$};
  \node[black,font=\small,align=center] at (3.2,6.0)
    {曲线在切线上方，弦在曲线上方 $\Longrightarrow$ $f$ 下凸};
\end{tikzpicture}
\end{center}

\noindent
图中：$x_0=\lambda x_1+(1-\lambda)x_2$ 为 $x_1,x_2$ 的加权点；
\emph{橙色切线}在切点外延伸到 $x_1,x_2$ 处的纵坐标\emph{不超过}曲线值（第 4 步的 (3)）；
\emph{红色弦}连接 $(x_1,f(x_1)),(x_2,f(x_2))$，其纵坐标\emph{不低于}曲线值（凸性定义）。
两条不等式正是证明的第 4、5 步。

\subsection*{图 2：斜率（$f'$）单调增加，即 $f''\ge0$}

\begin{center}
\begin{tikzpicture}[scale=1.05,>=Stealth]
  \draw[->,gray] (0,0)--(6.4,0) node[right,black]{$x$};
  \draw[->,gray] (0,0)--(0,4.8) node[above,black]{$f$};
  % f(x)=0.6+0.13*(x-1)^2，三处切点斜率递增
  % f'(x)=0.26*(x-1)
  \draw[curve,very thick,domain=0.6:5.9,samples=100,smooth]
       plot (\x,{0.6+0.13*(\x-1)*(\x-1)});
  \foreach \x/\l in {1.5/a,3.0/b,4.6/c}{
     \pgfmathsetmacro{\y}{0.6+0.13*(\x-1)*(\x-1)}
     \pgfmathsetmacro{\k}{0.26*(\x-1)}
     \draw[tang,thick] (\x-0.55,{\y-\k*0.55})--(\x+0.55,{\y+\k*0.55});
     \fill (\x,\y) circle (1.8pt) node[below right,black,font=\scriptsize]{$\l$};
  }
  \node[black,font=\small,align=center] at (3.2,5.35)
    {从左到右切线越来越陡：$f'$ 单调增 $\Longleftrightarrow$ $f''\ge0$};
\end{tikzpicture}
\end{center}

\noindent
从左到右三条切线越来越陡（$a\to b\to c$ 斜率递增），直观对应 $f'$ 单调增加、
等价于 $f''\ge0$。这正是第 1 步「$f''\ge0\Rightarrow f'$ 单调增」的几何图像，
也是第 4 步 $f'(\eta_1)\le f'(x_0)\le f'(\eta_2)$ 的来源。

\subsection*{图 3：证明用到的三点与两段}

\begin{center}
\begin{tikzpicture}[scale=1.15,>=Stealth]
  \draw[->,gray] (0,0)--(8.6,0) node[right,black]{$x$};
  \draw[->,gray] (0,0)--(0,4.4) node[above,black]{$f$};
  % 三点 x1=1, x0=4, x2=7 ; 上凸弧 y=0.5+0.13*(x-4)^2 + ... 用向上开口
  % y=1.0+0.13*(x-4)^2 ; x1: 1.0+0.13*9=2.17 ; x0:1.0 ; x2: 2.17
  \foreach \x/\l in {1.0/x_1,4.0/x_0,7.0/x_2}{
     \draw[gray] (\x,0.08)--(\x,-0.08);
     \fill (\x,0) circle (1.4pt) node[below,black,font=\small]{$\l$};
  }
  \draw[curve,very thick,domain=0.5:7.6,samples=90,smooth]
       plot (\x,{1.0+0.13*(\x-4.0)*(\x-4.0)});
  % x1,y=2.17 ; x0:1.0 ; x2:2.17
  \fill (1.0,2.17) circle (1.8pt);
  \fill (4.0,1.0) circle (1.8pt);
  \fill (7.0,2.17) circle (1.8pt);
  % 左段平均斜率线 (x1->x0) : (2.17-1.0)/(1-4)= 1.17/(-3) = -0.39
  \draw[secline,thick] (1.0,2.17)--(4.0,1.0)
       node[pos=0.5,below,font=\scriptsize,black]{$f'(\eta_1)$};
  % 右段平均斜率线 (x0->x2): (2.17-1.0)/(7-4)= 0.39
  \draw[auxc,thick] (4.0,1.0)--(7.0,2.17)
       node[pos=0.5,below,font=\scriptsize,black]{$f'(\eta_2)$};
  % 中心切线（x0=4 顶点，斜率0，水平）
  \draw[tang,thick] (3.1,1.0)--(4.9,1.0);
  \node[tang,font=\scriptsize,above] at (4.9,1.0){$f'(x_0)=0$};
  \node[black,font=\small,align=center] at (4.0,4.1)
    {$x_1<x_0<x_2$：两段各用一次拉格朗日中值定理\\ $f'(\eta_1)\le f'(x_0)\le f'(\eta_2)$};
\end{tikzpicture}
\end{center}

\noindent
在 $[x_1,x_0]$、$[x_0,x_2]$ 上分别取中值点 $\eta_1,\eta_2$，得平均斜率 $f'(\eta_1)$、$f'(\eta_2)$（图中两段直线）。
由 $f'$ 单调增，$f'(\eta_1)\le f'(x_0)\le f'(\eta_2)$——中心点切线斜率夹在两者之间。

% =====================================================================
\section*{附：三种等价刻画速查}
\begin{center}
\renewcommand{\arraystretch}{1.5}
\begin{tabular}{>{\centering\arraybackslash}p{3.6cm} >{\centering\arraybackslash}p{10.5cm}}
\toprule
\textbf{等价说法} & \textbf{表达式}\\
\midrule
凸性定义（弦） & $f(\lambda x_1+(1-\lambda)x_2)\le \lambda f(x_1)+(1-\lambda)f(x_2)$\\
切线支撑 & $f(x)\ge f(x_0)+f'(x_0)(x-x_0)$\\
斜率递增 & $f'$ 单调增加，即 $f''\ge 0$\\
\bottomrule
\end{tabular}
\end{center}

\end{document}
